\section{Possession by Evil Spirits}

In discussing the arcanum of the Empress, Tomberg brings up the idea of possession by evil spirits. Unlike what you see in Hollywood movies, this is not the result of an invasion by evil spirits; to the contrary, these spirits must be invited in. Once entrenched in consciousness, the can come to dominate the mind to their own ends. He writes:

\begin{quotationx}
The objective of scared Magic is more than healing, pure and simple. It is the restoration of freedom, including freedom from the stranglehold of doubt, fear, hate, apathy and despair. The ``evil spirits'' that deprive man of his freedom are not at all those of the so-called hierarchies of evil or fallen hierarchies. Neither Satan, nor Belial, nor Lucifer, nor Mephistopheles has ever deprived anyone of his freedom. Temptation is their only weapon and it presupposes the free will of the one tempted. But possession by an ``evil spirit'' has nothing to do with temptation. It is the same thing as with Frankenstein's monster. A man engenders an elemental and becomes thereby the slave of his own creation. The demons or evil spirits of the New Testament are today called in psychotherapy obsessive compulsive disorders, phobias, fixed ideas, and so on. They were discovered by contemporary psychiatrists and recognized as real, that is to say, parasitic psychic organisms independent of conscious human will and tending to control it.

But the Devil is not involved with that, at least as far as direct participation goes. He observes and never violates the law which safeguards human freedom and is the convention of both the hierarchies of the left and right. May you not fear the devil, but rather the perverse inclinations in yourself! For these perverse human inclinations can deprive us of our freedom and control us. Even worse, they can make use of our imagination and inventive capabilities and lead us to creations that can become the scourge of mankind.
\end{quotationx}


Hence, possession is the consequence of engendering elemental beings and then being enslaved by one's own creations. In our day, we would explain possession in psychological terms: obsessive/compulsive disorders, phobias, fixed ideas, for example. In a sense these neuroses become real and have an independent existence from the person who engendered them.

The elemental beings correspond to air, water, earth, fire, so each one would have a different effect on the mind. The elementals do not have independent conscious existence. They are more like a virus and only get ``life'' by taking over the soul or mind, actually by being ``engendered''. Being possessed by elementals is not a sign of spirituality; on the contrary, it is a serious impediment.

The way to overcome the elementals is to never yield to their characteristic defects. Instead, adopt their corresponding strengths. E.g., be prompt and active like the Sylphs, energetic and strong like the Salamanders, and so on.

The elementals are spirits in bondage to the elements. As such, they are not free and cannot have the use of reason; the can only manifest in subhuman forms as animals. Nevertheless, as Tomberg points out, they can become invited into a human consciousness, thus effectively dominating it, creating ``vicious and imperfect men'' (Eliphas Levi). Levi writes:

\begin{quotex}
Elementary spirits are like children: they torment chiefly whose who trouble about them ... they frequently occasion our bizarre or disturbing dreams and produce psychic phenomena . They can manifest no thought other than our own and when we are not thinking they speak to us with the incoherence of dreams. They reproduce good and evil indifferently, for they are without free will and are hence irresponsible... Such creatures are neither damned nor guilty; they are curious and innocent.
\end{quotex}


That is what makes them so tempting: they are innocent and can appear to be beautiful as do the sylphs and undines. Their child-like and animal-like qualities make them endearing. The myths and legends of these spirits typically use trickery or enticement to invade one's consciousness.

Chart \ref{tab:QualitiesSpirits}, adapted from \emph{Transcendental Magic}, lists the qualities of the four elemental spirits.

\begin{table}[h]\small\centering
\begin{tabular}{p{.1\textwidth}p{.15\textwidth}*{3}{p{.2\textwidth}}}\bottomrule
\textbf{Element} & \textbf{Elemental} & \textbf{Temperament} & \textbf{Defect} & \textbf{Strength}\\\midrule
Fire & Salamander & Sanguine & Passionate & Energetic and strong\\
Air & Sylph & Bilious & Shallow and capricious & Prompt and active\\
Water & Undine & Phlegmatic & Irresolute, cold and fickle & Pliant and attentive to images\\
Earth & Gnome & Melancholy & Avaricious and greedy & Laborious and patient\\\bottomrule
\end{tabular}
\label{tab:QualitiesSpirits}
\caption{Qualities of the Four Spirits}
\end{table}

Originally published on 26 February 2011 at Mediations on the Tarot.


\flrightit{Posted on 2022-12-12 by Cologero}

